\section{Richardson--Romberg extrapolation}

To eliminate the leading $O(\alpha)$ bias term, we employ the
\emph{Richardson--Romberg} (RR) \emph{extrapolation} procedure. Two LSA
sequences are run \emph{on the same Markov chain trajectory} $\{Z_k\}$ with
step sizes $\alpha$ and $2\alpha$, and the RR iterate is formed as
\begin{equation}
  \overline{\theta}_n^{(\alpha,\mathrm{RR})}
  =2\overline{\theta}_n^{(\alpha)}
   -\overline{\theta}_n^{(2\alpha)}.
  \label{eq:rr-iterate}
\end{equation}
Since both sequences share the same noise realization, the leading bias term
$\alpha\Delta$ cancels, leaving a residual bias of order
$O(\alpha^{3/2})$ or higher (Levin et al., 2025).

More generally, one can consider the multi-level extrapolation with $M$ step
sizes $\mathcal{A}=\{\alpha_1,\ldots,\alpha_M\}$ and coefficients
$\{h_m\}$ determined by the Vandermonde system (Huo et al., 2024):
\[
  \sum_{m=1}^M h_m=1,
  \qquad
  \sum_{m=1}^M h_m\alpha_m^l=0,
  \qquad l=1,\ldots,M-1,
\]
which cancels successive powers in settings where such a power-series
expansion is available.
